\begin{CJK}{UTF8}{gbsn}
\chapter{摘要}
%\addcontentsline{toc}{chapter}{Introduction}  
该论文研究了在质子-质子对撞中，质心系能量分别为7,8以及13TeV下，通过四轻子衰变 $ H \to 4\ell$ ($\ell = \Pe$, $\mu$) 的希格斯玻色子产生截面的测量问题。
希格斯玻色子的基准产生截面测量数据都源自于大型强子对撞机LHC上运行的CMS探测器，对应于不同质心系能量，积分亮度分别为7TeV下5.1~\fbinv, 8TeV下19.7~\fbinv, 以及13TeV下
2.8~\fbinv。这些测量结果为标准模型在TeV能级下的希格斯物理相关的QCD微扰计算提供了突破口。测量结果与标准模型的理论预期有任何的偏离都将会为超出标准模型的新物理提供直接的线索。
\par 
质心系能量为8TeV时的基准微分截面已给出，对希格斯玻色子产生机制敏感的几个动力学观测量进行了考察，包括快度，四轻子体系的横动量，相关的多重喷注，领头喷注的横动量，以及对领头喷注和希格斯玻色
子候选粒子间的分离。对希格斯玻色子衰变机制敏感的几个动力学观测量的研究结果也已公布~\cite{Reference40}，这些动力学观测量包括：两组轻子对的不变质量，以及Collins-Soper体系下轻子形成的五个螺旋度角。在希
格斯的测量基础上，本文也利用了同样的质心系能量为8TeV的数据和微分变量，来对 $Z \rightarrow 4\ell$ 的基准微分截面进行测量。 $Z \rightarrow 4\ell$的基准微分截面的测量对检验标准模型的
理论预言以及对希格斯玻色子产生截面的测量方法的合理性的验证都至关重要。
\par $ H \to 4\ell$ 和 $Z \rightarrow 4\ell$ 的积分基准截面比率的计算结果也已发表。所有的截面测量都取自于基准相空间，其定义是来自于轻子运动学和事例拓扑结构的筛选的结果。总的基
准截面测量结果在7, 8和13TeV下分别为$0.56^{\rm +0.67}_{\rm -0.44}({\rm stat.})^{\rm +0.21}_{\rm -0.06}({\rm sys.})^{+0.02}_{-0.02}({\rm model})$ ${\rm fb}$,
$1.11^{\rm +0.41}_{\rm -0.35}({\rm stat.)^{\rm +0.14}_{\rm -0.10}({\rm sys.}) ^{+0.08}_{-0.02}({\rm model}){\rm fb}$
以及 $2.48^{+1.45}_{-1.12}({\rm stat.})^{+0.28}_{-0.17}({\rm sys.})^{+0.01}_{-0.04}({\rm model})~{\rm fb}$。所有的截面测量结果和基于标准模型是理论预期之间没有显著性的偏差。
\end{CJK}